% GENERATED FILE. Edit canonical lessons, then run npm run generate:books.
% canonical-source-hash: fnv1a64:68209c482c59f9bc
% canonical-lessons: FA-C164-behtar, FA-C164-movaffaq, FA-C164-kond, FA-C164-zebr, FA-C164-khalvat

\chapter{Better, Successful, Blunt, Rough, Uncrowded}
\label{ch:fa-better-successful-blunt-rough-uncrowded}
\hlchaptermodality{164}

\begin{chapteropening}
\textbf{By the end of this chapter:} \emph{I can say better, successful, blunt, rough and uncrowded.}

Complete the last lesson of chapter 164: I can say better, successful, blunt, rough and uncrowded.
\end{chapteropening}

\section[\texorpdfstring{\fa{بهتر} (behtar)}{بهتر (behtar)}]{\fa{بهتر} (behtar) — better}

\label{lesson:FA-C164-behtar}



\begin{quote}
Before the new one: say the Persian for broken, then the Persian for different.
\end{quote}

\subsection*{You'll want to know: \fa{بهتر}}
\textbf{\fa{بهتر}} — \emph{behtar} — \textquotedblleft{}better\textquotedblright{}.

\begin{cousinweb}
\emph{behtar} is \emph{beh}, \textquotedblleft{}good\textquotedblright{}, with the comparative ending \emph{-tar}. \emph{beh} is native, usually traced to Middle Persian \emph{weh}. English \emph{better} only looks alike: it is a Germanic word, and no source links it to \emph{beh}.
\end{cousinweb}

\begin{grammarlens}[title={What you've built}]
One more word for this chapter.
\end{grammarlens}

\subsection*{Your turn}
\begin{itemize}
\raggedright
  \item \emph{Say it:} \emph{behtar}
  \item \emph{Say it:} \emph{behtar}, once more
  \item \emph{Say it:} say \emph{behtar}
  \item \emph{Recall:} say the Persian for broken, then the Persian for different, then say \emph{behtar} again
\end{itemize}

\begin{tcolorbox}[breakable,colback=teal!4,colframe=teal!35!black,title={Before you move on}]
What does \fa{بهتر} mean? (\textquotedblleft{}Better\textquotedblright{}.) Say it once more.
\end{tcolorbox}
\section[\texorpdfstring{\fa{موفق} (movaffaq)}{موفق (movaffaq)}]{\fa{موفق} (movaffaq) — successful}

\label{lesson:FA-C164-movaffaq}



\begin{quote}
Before the new one: say the Persian for different, then the Persian for better.
\end{quote}

\subsection*{You'll want to know: \fa{موفق}}
\textbf{\fa{موفق}} — \emph{movaffaq} — \textquotedblleft{}successful\textquotedblright{}.

\begin{cousinweb}
\emph{movaffaq} is borrowed from Arabic \emph{muwaffaq}, \textquotedblleft{}successful, favoured\textquotedblright{}. Its Arabic root is \emph{w-f-q} (Lane).
\end{cousinweb}

\begin{grammarlens}[title={What you've built}]
One more word for this chapter.
\end{grammarlens}

\subsection*{Your turn}
\begin{itemize}
\raggedright
  \item \emph{Say it:} \emph{movaffaq}
  \item \emph{Say it:} \emph{movaffaq}, once more
  \item \emph{Say it:} say \emph{movaffaq}
  \item \emph{Recall:} say the Persian for different, then the Persian for better, then say \emph{movaffaq} again
\end{itemize}

\begin{tcolorbox}[breakable,colback=teal!4,colframe=teal!35!black,title={Before you move on}]
What does \fa{موفق} mean? (\textquotedblleft{}Successful\textquotedblright{}.) Say it once more.
\end{tcolorbox}
\section[\texorpdfstring{\fa{کند} (kond)}{کند (kond)}]{\fa{کند} (kond) — blunt; slow}

\label{lesson:FA-C164-kond}



\begin{quote}
Before the new one: say the Persian for better, then the Persian for successful.
\end{quote}

\subsection*{You'll want to know: \fa{کند}}
\textbf{\fa{کند}} — \emph{kond} — \textquotedblleft{}blunt; slow\textquotedblright{}.

\begin{cousinweb}
\emph{kond} is a native Persian word, inherited rather than borrowed: EtymDB takes it back to an Indo-Iranian word for \textquotedblleft{}blunt, dull\textquotedblright{}. (It also lists \emph{qand}, \textquotedblleft{}sugar candy\textquotedblright{}, a different word.)
\end{cousinweb}

\begin{grammarlens}[title={What you've built}]
One more word for this chapter.
\end{grammarlens}

\subsection*{Your turn}
\begin{itemize}
\raggedright
  \item \emph{Say it:} \emph{kond}
  \item \emph{Say it:} \emph{kond}, once more
  \item \emph{Say it:} say \emph{kond}
  \item \emph{Recall:} say the Persian for better, then the Persian for successful, then say \emph{kond} again
\end{itemize}

\begin{tcolorbox}[breakable,colback=teal!4,colframe=teal!35!black,title={Before you move on}]
What does \fa{کند} mean? (\textquotedblleft{}Blunt; slow\textquotedblright{}.) Say it once more.
\end{tcolorbox}
\section[\texorpdfstring{\fa{زبر} (zebr)}{زبر (zebr)}]{\fa{زبر} (zebr) — rough}

\label{lesson:FA-C164-zebr}



\begin{quote}
Before the new one: say the Persian for successful, then the Persian for blunt.
\end{quote}

\subsection*{You'll want to know: \fa{زبر}}
\textbf{\fa{زبر}} — \emph{zebr} — \textquotedblleft{}rough\textquotedblright{}.

\begin{cousinweb}
\emph{zebr} is a native Persian word, but no source used here gives its older form, and its origin is not settled.
\end{cousinweb}

\begin{grammarlens}[title={What you've built}]
One more word for this chapter.
\end{grammarlens}

\subsection*{Your turn}
\begin{itemize}
\raggedright
  \item \emph{Say it:} \emph{zebr}
  \item \emph{Say it:} \emph{zebr}, once more
  \item \emph{Say it:} say \emph{zebr}
  \item \emph{Recall:} say the Persian for successful, then the Persian for blunt, then say \emph{zebr} again
\end{itemize}

\begin{tcolorbox}[breakable,colback=teal!4,colframe=teal!35!black,title={Before you move on}]
What does \fa{زبر} mean? (\textquotedblleft{}Rough\textquotedblright{}.) Say it once more.
\end{tcolorbox}
\section[\texorpdfstring{\fa{خلوت} (khalvat)}{خلوت (khalvat)}]{\fa{خلوت} (khalvat) — uncrowded}

\label{lesson:FA-C164-khalvat}



\begin{quote}
Before the new one: say the Persian for blunt, then the Persian for rough.
\end{quote}

\subsection*{You'll want to know: \fa{خلوت}}
\textbf{\fa{خلوت}} — \emph{khalvat} — \textquotedblleft{}uncrowded\textquotedblright{}.

\begin{cousinweb}
\emph{khalvat} is borrowed from Arabic \emph{ḫalwa}, \textquotedblleft{}solitude, privacy\textquotedblright{}. Its Arabic root is \emph{ḫ-l-w} (Lane). Persian uses it for an uncrowded place.
\end{cousinweb}

\begin{grammarlens}[title={What you've built}]
One more word for this chapter.
\end{grammarlens}

\subsection*{Your turn}
\begin{itemize}
\raggedright
  \item \emph{Say it:} \emph{khalvat}
  \item \emph{Say it:} \emph{khalvat}, once more
  \item \emph{Say it:} say \emph{khalvat}
  \item \emph{Recall:} say the Persian for blunt, then the Persian for rough, then say \emph{khalvat} again
\end{itemize}

\begin{tcolorbox}[breakable,colback=teal!4,colframe=teal!35!black,title={Before you move on}]
What does \fa{خلوت} mean? (\textquotedblleft{}Uncrowded\textquotedblright{}.) Say it once more.
\end{tcolorbox}
